\section{The Moving-Knife Algorithm}
\label{sec:algorithm}

\subsection{Notation}

Let $G_v = (V_v, E_v)$ be the subgraph at bisection tree node $v$, with
$m = |V_v|$ tracts and population $p_i$ for each tract $i \in V_v$.
Let $\ell_i = (\lambda_i, \phi_i)$ be the WGS84 (longitude, latitude) of
the population-weighted centroid of tract $i$'s polygon, loaded from the
centroid companion file (Section~\ref{subsec:data}).
Let $\mathbf{x}_i = (x_i, y_i) = \mathrm{proj}(\ell_i)$ be the
Albers-projected coordinates in metres (EPSG:5070, shared with
\cvdgeo{}~\citep{dellaLibera2026cvdgeo}).

\subsection{Design Choices}
\label{subsec:design}

Three non-obvious design decisions underlie the algorithm.

\paragraph{Choice 1: $n_{\mathrm{orient}} = 180$ (every $1^\circ$).}
The sweep tests orientations $\theta \in \{0, 1, \ldots, 179\}$ degrees,
covering $[0^\circ, 180^\circ)$ uniformly.
Restricting to $[0^\circ, 180^\circ)$ is correct because opposite
orientations produce the same bisecting hyperplane with left and right
swapped; no information is lost by this restriction.
The 1$^\circ$ granularity ensures the globally best direction is never
missed by more than half a degree.
Runtime is $O(n_{\mathrm{orient}} \cdot m \log m) \approx O(180 m \log m)$ per bisection node,
which is comparable to a single iteration of \cvdgeo{} and substantially
faster than the full \cvdgeo{} run (typically 7--12 iterations $\times$ $O(m)$).

\paragraph{Choice 2: Reock as the default metric.}
The sweep scores each split by $\min(\reock{}_{\mathrm{left}}, \reock{}_{\mathrm{right}})$:
the worst-case Reock score of the two halves.
Using the minimum (rather than the mean) ensures both halves are compact,
following the maximin principle: maximise the outcome for the worst-off
district.
Reock is the default (rather than Polsby-Popper) because Reock requires only
centroid positions and MEC geometry; Polsby-Popper requires polygon perimeter
data, introducing an additional data dependency.
The flag \texttt{--mka-metric polsby} enables Polsby-Popper scoring for
jurisdictions where that metric is legally operative.

\paragraph{Choice 3: Post-hoc rebalance.}
The linear sweep produces a population split that is close to but not exactly
balanced: the binary search for the split point may overshoot by at most one
tract's population.
The same 200-iteration boundary-swap rebalance used by \cvdgeo{} and \bfsg{}
enforces exact balance within tolerance.
The rebalance also repairs contiguity: for non-convex subgraph geometries, a
linear sweep may produce disconnected half-subgraphs; BFS-based contiguity
repair is applied before returning the final split.

\subsection{Seed Derivation}
\label{subsec:seed}

\mka{} is deterministic given \texttt{base\_seed} and $n_{\mathrm{orient}}$.
A seed is needed only for tie-breaking when two orientations produce
identical minimum-Reock scores (rare in practice).

\begin{definition}[\mka{} Node Seed]
\label{def:mka-seed}
Let $\mathtt{base\_seed}$ be the unsigned 64-bit global seed and
$\mathtt{node\_path}$ the path string for bisection tree node $v$.
The \emph{\mka{} node seed} is:
\[
  s_v^{\mathrm{mka}} = \mathrm{first8}\bigl(
    \mathrm{SHA\text{-}256}(
      \texttt{"MKA\_INIT\_"}
      \;\|\; \mathrm{le4}(|\mathtt{node\_path}|)
      \;\|\; \mathtt{node\_path}
      \;\|\; \mathrm{le8}(\mathtt{base\_seed})
    )
  \bigr),
\]
where $\mathrm{le4}(n)$ encodes $n$ as a 4-byte little-endian \texttt{u32},
$\mathrm{le8}(x)$ encodes $x$ as an 8-byte little-endian \texttt{u64}, and
$\mathrm{first8}(h)$ reads the first 8 bytes of the SHA-256 digest as a
little-endian \texttt{u64}~\citep{sha256nist}.
\end{definition}

The length prefix $\mathrm{le4}(|\mathtt{node\_path}|)$ eliminates
prefix-collision risk: the boundary between \texttt{node\_path} and
\texttt{base\_seed} is unambiguous.
The prefix \texttt{"MKA\_INIT\_"} is distinct from \texttt{"CVD\_GEO\_INIT\_"}
(Phase~2 \cvdgeo{}) and \texttt{"CVD\_INIT\_"} (Phase~1 \cvd{}).
A platform invariant test asserts all three prefix constants are present in
source and are pairwise distinct.

\subsection{Formal Pseudocode}

Algorithm~\ref{alg:mka} gives the complete \mka{} bisection procedure.

\begin{algorithm}[t]
\caption{\mka{}-\textsc{Bisect}($G_v$, $\mathbf{x}$, $\mathbf{p}$,
  $A$, $n_{\mathrm{orient}}$, $\mathtt{base\_seed}$,
  $\mathtt{node\_path}$, $\mathtt{balance\_tolerance}$)}
\label{alg:mka}
\begin{algorithmic}[1]
\State $m \gets |V_v|$;\quad $P \gets \sum_{i \in V_v} p_i$
\If{$m < 2$}
  \textbf{return} degenerate assignment
\EndIf
\State $s_v \gets \mka{}\text{-Seed}(\mathtt{base\_seed}, \mathtt{node\_path})$
       \Comment{Definition~\ref{def:mka-seed}}
\State $\mathtt{sorted} \gets \mathrm{sort}(V_v)$
       \Comment{sort by global tract index}
\State $\mathbf{x}' \gets [\mathbf{x}_i : i \in \mathtt{sorted}]$
       \Comment{projected coordinates in sorted order}
\State $\mathtt{best\_score} \gets -\infty$;\quad
       $\mathtt{best\_left} \gets \emptyset$;\quad
       $\mathtt{best\_right} \gets \emptyset$
\For{$k = 0$ \textbf{to} $n_{\mathrm{orient}} - 1$}
  \State $\theta \gets k \cdot \pi / n_{\mathrm{orient}}$
         \Comment{angle in $[0, \pi)$}
  \State $\mathbf{u} \gets (\cos\theta, \sin\theta)$
         \Comment{unit sweep vector}
  \State $\mathrm{proj}_i \gets \mathbf{x}'_i \cdot \mathbf{u}$ for all $i$
         \Comment{project onto sweep axis}
  \State $\sigma \gets$ sort indices by $\mathrm{proj}_i$ ascending
  \State $C \gets 0$;\quad $\mathtt{split} \gets 0$
  \For{$j = 0$ \textbf{to} $m-1$}
    \State $C \gets C + p_{\sigma[j]}$
    \If{$C \ge P/2$}
      $\mathtt{split} \gets j$;\quad \textbf{break}
    \EndIf
  \EndFor
  \State $L \gets \{\mathtt{sorted}[\sigma[j]] : j \le \mathtt{split}\}$
  \State $R \gets \{\mathtt{sorted}[\sigma[j]] : j > \mathtt{split}\}$
  \State $r_L \gets \reock{}(L, \mathbf{x}, A)$
         \Comment{Reock via Welzl MEC; clamped to $[0,1]$}
  \State $r_R \gets \reock{}(R, \mathbf{x}, A)$
  \State $\mathtt{score} \gets \min(r_L, r_R)$
  \If{$\mathtt{score} > \mathtt{best\_score}$}
    \State $\mathtt{best\_score} \gets \mathtt{score}$;\quad
           $\mathtt{best\_left} \gets L$;\quad
           $\mathtt{best\_right} \gets R$
  \ElsIf{$\mathtt{score} = \mathtt{best\_score}$}
    \State tie-break by $s_v$
  \EndIf
\EndFor
\State $(\mathtt{best\_left}, \mathtt{best\_right}) \gets
       \mathrm{Rebalance}(\mathtt{best\_left}, \mathtt{best\_right},
       G_v, \mathbf{p}, \mathtt{balance\_tolerance})$
       \Comment{200-iter boundary-swap}
\State \textbf{return} $(\mathtt{best\_left}, \mathtt{best\_right})$
\end{algorithmic}
\end{algorithm}

The \textbf{Reock} subroutine (line 17) computes:
\[
  \reock{}(S, \mathbf{x}, A) =
  \min\!\left(1.0,\;
    \frac{\sum_{i \in S} A_i}{\pi \cdot r_{\mathrm{MEC}}(S, \mathbf{x})^2}
  \right),
\]
where $A_i$ is the TIGER/Line polygon area of tract $i$ in square metres,
and $r_{\mathrm{MEC}}(S, \mathbf{x})$ is the radius of the smallest enclosing
circle of the projected centroid points $\{\mathbf{x}_i : i \in S\}$,
computed by Welzl's algorithm (Proposition~\ref{prop:welzl}).
The clamp enforces the $\reock{} \le 1$ invariant under the centroid
approximation (Remark~\ref{rem:centroid-approx}).

\subsection{Runtime Complexity}

\begin{proposition}[\mka{} Runtime]
\label{prop:mka-runtime}
\mka{} bisection at a node with $m$ tracts requires:
\begin{itemize}
\item Per orientation ($n_{\mathrm{orient}}$ total):
  \begin{itemize}
  \item Projection: $O(m)$ dot products.
  \item Sort: $O(m \log m)$ (dominated by projection if $m$ is small;
    note that the sort can be replaced by a linear scan when $n_{\mathrm{orient}}$
    is large, since successive orientations produce nearly-sorted projections).
  \item Population binary search: $O(m)$ cumulative sum.
  \item Welzl MEC (two calls): expected $O(m)$ each.
  \end{itemize}
\item Post-hoc rebalance: $O(200m)$.
\item Total: $O(n_{\mathrm{orient}} \cdot m \log m + 200m)
            = O(n_{\mathrm{orient}} \cdot m \log m)$.
\end{itemize}
For $n_{\mathrm{orient}} = 180$: $O(180 m \log m)$.
This is comparable to one full \cvdgeo{} iteration count (typically 7--12)
$\times$ $O(m)$ per iteration $= O(84m)$--$O(144m)$.
The sort step dominates; using a radix sort reduces the per-orientation
cost to $O(m)$, giving total $O(180m)$ --- strictly faster than \cvdgeo{}.
\end{proposition}

\subsection{Data Requirements}
\label{subsec:data}

\mka{} requires the same centroid companion file introduced by \cvdgeo{}:
\texttt{\{state\}\_centroids\_\{year\}.json}, stored alongside the adjacency
file in \texttt{data/\{year\}/}.
It additionally requires the tract area field (\texttt{LoadedGraph.tract\_areas}),
which stores TIGER/Line polygon areas in square metres and is bundled with
the existing adjacency file.
No additional fetch step is needed if centroids were already fetched for
\cvdgeo{} Phase~2.

If the centroid companion file is absent and \texttt{--structure moving-knife}
is requested, the CLI returns a diagnostic error:
\begin{verbatim}
ERROR: --structure moving-knife requires tract centroid data.
       Run: bisect fetch --type centroids --states NC --year 2020
\end{verbatim}
No silent fallback to a different algorithm occurs.

\subsection{CLI and YAML}
\label{subsec:cli}

\paragraph{Pure \mka{} structure:}
\begin{verbatim}
bisect state --state NC --year 2020 \
  --structure moving-knife \
  --mka-orientations 180 \
  --mka-metric reock
\end{verbatim}

\paragraph{Hybrid: \mka{} direction + \areasec{} balance enforcement:}
The \texttt{split\_subgraph\_mka\_direction} function returns the optimal
angle $\theta^* \in [0, \pi)$ without performing the full split.
\areasec{} uses $\theta^*$ to override its \texttt{ratio\_direction} before
applying its own balance logic:
\begin{verbatim}
bisect state --state NC --year 2020 \
  --structure area-section \
  --area-section-init moving-knife
\end{verbatim}

\paragraph{YAML configuration:}
\begin{verbatim}
algorithm:
  structure: moving-knife
  mka_orientations: 180
  mka_metric: reock           # or: polsby-popper
  weights: geographic
  search: convergence
  convergence_threshold: 600
  balance_tolerance: 0.5
workers: 8
\end{verbatim}

\subsection{Audit Chain}
\label{subsec:audit}

Every \mka{} run appends the following fields to
\texttt{runs/\{label\}/\{year\}/index.json}:

\begin{verbatim}
"structure":         "moving-knife",
"mka_orientations":  180,
"mka_metric":        "reock",
"mka_seed_formula":  "SHA-256('MKA_INIT_'||node_path_len:u32le
                       ||node_path||base_seed:u64le)",
"node_path":         "01",
"base_seed":         12345678,
"mka_seed":          987654321,
"optimal_angle_deg": 47.0,
"reock_left":        0.412,
"reock_right":       0.389,
"min_reock_score":   0.389,
"convergence_warning": false
\end{verbatim}

\texttt{optimal\_angle\_deg} enables independent verification: an auditor
reruns the sweep with the same centroids and $n_{\mathrm{orient}}$ and confirms
$\theta^*$ matches.
\texttt{node\_path} is required for independent seed reproduction.
\texttt{convergence\_warning} is always \texttt{false} for \mka{}:
the sweep is exhaustive over all orientations and has no iterative
convergence criterion.
\texttt{min\_reock\_score} is the quantity \mka{} directly maximises.
